\begin{abstract}
Sự bùng nổ của các hệ thống IoT với hàng trăm nghìn cảm biến tạo ra dòng dữ liệu liên tục ở tốc độ cao, đặt ra yêu cầu xử lý đồng thời về \textit{khối lượng}, \textit{tốc độ} và \textit{độ tin cậy} (veracity) --- ba trong số các đặc trưng cốt lõi của Big Data. Báo cáo này trình bày việc thiết kế và đánh giá một nền tảng Big Data cho bài toán streaming dữ liệu cảm biến IoT trên một máy Linux duy nhất, dựa trên các thành phần mã nguồn mở phổ biến: giao thức \texttt{MQTT} qua broker \texttt{Mosquitto} làm lớp tiếp nhận thiết bị, một \texttt{Bridge} Python chuyển đổi MQTT sang \texttt{Apache Kafka} (KRaft) làm bộ đệm phân tán có thứ tự, \texttt{Spark Structured Streaming} ở chế độ local làm stream processor theo mô hình Bronze--Silver--Gold trên Parquet cục bộ, \texttt{PostgreSQL} làm tầng phục vụ và lưu aggregate Gold, \texttt{FastAPI} cùng dashboard tĩnh làm lớp truy vấn; tác vụ tổng hợp lại theo giờ được kích hoạt bằng script. Toàn bộ stack được đóng gói bằng Docker Compose với ngân sách bộ nhớ khoảng 4,6 GiB cho chế độ streaming.

Dữ liệu cảm biến được mô hình hoá dưới dạng vector đa chỉ số môi trường hợp nhất: mỗi bản tin mang đồng thời sáu đại lượng --- nhiệt độ, độ ẩm, CO\textsubscript{2}, áp suất khí quyển, bụi mịn PM2.5 và ánh sáng --- kèm định danh sự kiện, định danh cảm biến và hai mốc thời gian (thời điểm hiện tượng và thời điểm tiếp nhận). Việc dùng nhiều loại chỉ số với miền giá trị, đơn vị và hướng ngưỡng cảnh báo khác nhau buộc các cơ chế kiểm tra chất lượng lẫn phát hiện cảnh báo phải tổng quát theo cấu hình thay vì viết cứng cho một đại lượng đơn lẻ, qua đó tránh nhầm lẫn giữa một pipeline ``tình cờ chạy đúng'' và một kiến trúc xử lý đúng đắn. Luồng xử lý tách bạch dữ liệu thô (Bronze, ghi nguyên bản để làm bản sao lưu bền vững khi Kafka hết hạn lưu), dữ liệu đã làm sạch (Silver) và dữ liệu lỗi bị cách ly kèm lý do máy đọc được (Quarantine).

Nghiên cứu tập trung vào các câu hỏi về hệ thống thay vì bản thân dữ liệu cảm biến: thông lượng và độ trễ tiếp nhận thay đổi thế nào khi tốc độ đầu vào và số nguồn phát tăng; ảnh hưởng của số partition Kafka và số nhân xử lý Spark đến khả năng mở rộng; hiệu quả phục hồi sau gián đoạn nhờ checkpoint và vai trò lớp đệm của Kafka; các cơ chế veracity theo từng chỉ số (kiểm tra trường bắt buộc và miền vật lý, suy luận lý do lỗi có cấu trúc, xử lý dữ liệu đến muộn, kế toán bản ghi bảo toàn) xử lý dữ liệu hỏng, trùng lặp và đến trễ ra sao; và tác động của việc tổng hợp lại theo giờ đến tính xác định của kết quả. Trên nền kiến trúc đó, báo cáo vận dụng khung khái niệm xử lý event-time có nhận thức veracity --- phân loại thời gian, observed lateness, ngữ nghĩa cửa sổ/watermark, vector veracity và metric chất lượng theo cửa sổ trượt --- làm ngôn ngữ chung cho thiết kế và đánh giá, đồng thời tham chiếu các thực hành benchmark và giám sát chất lượng dữ liệu stream-first trong tài liệu gần đây.

Báo cáo đề xuất một quy trình benchmark gồm các kịch bản chạy được trên một máy: đa nguồn phát đồng thời và cân bằng partition, tích luỹ và tiêu tán backlog khi dừng/bật lại bộ xử lý, tăng tải để xác định điểm bão hoà và tầng nghẽn, suy rộng số partition và số bridge, cùng tiêm lỗi veracity theo từng chỉ số. Một phép chạy dài đã được thực hiện với năm nguồn, kéo dài bốn giờ, ghi nhận 50.760.004 sự kiện và 16.226.141.856 byte payload tại Bronze; kết quả này xác nhận mốc 15 GB trong bốn giờ của cấu hình một máy, nhưng không được suy diễn thành mục tiêu 15 GB trong ba giờ. Mục tiêu tải tiếp theo của hệ thống là 30 GB trong bốn giờ và chưa được đo trong báo cáo này. Kiến trúc cố ý tách bạch cái demo được trên một máy khỏi cái đòi hỏi hạ tầng nhiều máy: bảo mật truyền thông, sẵn sàng cao và cụm phân tán thật được ghi nhận là hướng phát triển chứ không triển khai trong bản này. Các số liệu đo thật được tách khỏi các bảng tham chiếu dựa trên tài liệu công bố và đi kèm đường dẫn evidence.
\end{abstract}
